%% toc: P2 20 20.5,20.6
\section{Set-theoretic prerequisites of the L\"owenheim--Skolem theorems}
\label{sec:p02-w02b-p02-029-ls-prerequisites}

The proofs of downward and upward L\"owenheim--Skolem
(\cref{thm:downward-ls,thm:upward-ls}) are short because they hand all of
their set theory to results stated in \cref{sec:set-prerequisites-06}: the
well-ordering theorem, cardinal assignment (\cref{cor:cardinal-assignment}),
the finite-coding theorem (\cref{thm:infinite-cardinal-finite-coding}) and
syntax counting (\cref{cor:syntax-counting}). Where they are stated, the
first three carry only a citation. This section does two things. It records
where each of those results is proved in full in this book. And it proves
every further set-theoretic step that the two L\"owenheim--Skolem proofs
take without comment: the closure properties of the bound
\(\abs X\leq\mu\), the existence of the cardinal
\(\mu=\max(\abs A,\abs\Lang,\aleph_0)\), the choice of all witness
functions from a single well-order, the bound on every stage of the
construction, the fresh constants and the counting in the upward proof, and
the fact that none of the three terms of the maximum can be removed.

\emph{Inputs.} The Axiom of Choice in the form of
\cref{sec:set-prerequisites-06}; comparison of sets by injections and
bijections and the Cantor--Schr\"oder--Bernstein theorem of
\cref{sec:set-prerequisites-05}; ordinals and cardinals
(\cref{def:ordinal-cardinal}); Kuratowski pairs
(\cref{sec:set-prerequisites-02}); signatures, the cardinality
\(\abs\Lang=\abs{\mathcal F\cup\mathcal R}\) and the countable set of
variables (\cref{sec:fol-syntax-semantics-01}); formulas
(\cref{sec:fol-syntax-semantics-02}); structures, which have nonempty domains
(\cref{sec:fol-semantics-01}); elementary substructures and the
Tarski--Vaught test (\cref{thm:tarski-vaught-fol}). The domain of a
structure \(\mathcal M\) is written \(\dom\mathcal M\); the inherited
statements write it \(\abs{\mathcal M}\), and \(\abs{\abs{\mathcal N}}\)
there is \(\abs{\dom\mathcal N}\) here. For sets \(X,Y\), \(\abs X\leq\abs Y\)
means that an injection \(X\to Y\) exists; for an ordinal \(\mu\),
\(\abs X\leq\mu\) means that an injection \(X\to\mu\) exists
(\cref{def:p01-w02-p01-011-readings}).

\emph{Conclusions.} Under the Axiom of Choice, every set-theoretic claim made
in the proofs of \cref{thm:downward-ls,thm:upward-ls} holds, with a proof in
this book: \cref{sec:p02-w02b-p02-029-where} for the four cited results, and
\cref{thm:p02-w02b-p02-029-bounds,thm:p02-w02b-p02-029-mu,thm:p02-w02b-p02-029-finite-injections,thm:p02-w02b-p02-029-skolem-functions,thm:p02-w02b-p02-029-stages,thm:p02-w02b-p02-029-fresh,thm:p02-w02b-p02-029-upward-count}
for the steps in between. \Cref{thm:p02-w02b-p02-029-terms-needed} shows that
the bound of downward L\"owenheim--Skolem cannot be improved by deleting a
term of the maximum.

\subsection{Where the cited results are proved}
\label{sec:p02-w02b-p02-029-where}

The proofs of \cref{thm:downward-ls,thm:upward-ls} use the following
results from \cref{sec:set-prerequisites-06}. Each is proved in Part~1 with
the same hypotheses; the proofs there are complete and are not repeated
here.
\begin{enumerate}
  \item \emph{Well-ordering theorem.} Under the Axiom of Choice every set
        has a strict well-order. Proved as
        \cref{thm:p01-w02-p01-008-well-ordering}, by a tower of subsets
        built from one choice function on the nonempty subsets of the set;
        no ordinals or recursion are used.
  \item \emph{Cardinal assignment} (\cref{cor:cardinal-assignment}). Under
        the Axiom of Choice every set \(A\) is bijective with exactly one
        cardinal \(\abs A\); the order of cardinals agrees with comparison by
        injections; any two cardinals are comparable; and every nonempty
        finite list of cardinals has a greatest member. Proved as items~1--4
        of \cref{thm:p01-w02-p01-008-cardinal-assignment}, using the
        well-ordering theorem and the representation of well-orders by
        ordinals (\cref{thm:p01-w02-p01-008-representation}).
  \item \emph{Finite coding, clauses~1, 2 and~4}
        (\cref{thm:infinite-cardinal-finite-coding}). For an infinite
        cardinal \(\kappa\): \(\abs{\kappa^n}=\kappa\) for every positive
        \(n\in\N\); \(\abs{\kappa^{<\omega}}=\kappa\); and
        \(\abs{\bigcup_{n\in\N}A_n}\leq\kappa\) whenever
        \(\abs{A_n}\leq\kappa\) for every \(n\). Proved as clauses~1, 2
        and~4 of \cref{thm:p01-w02-p01-011-finite-coding}. Clause~1 rests on
        the choice-free identity \(\abs{\mu\times\mu}=\abs\mu\) for infinite
        cardinals \(\mu\) (\cref{thm:p01-w02-p01-011-square}); clauses~2
        and~4 each apply the Axiom of Choice once, to select a whole
        sequence of injections. Clause~3 of the theorem is not used by
        either L\"owenheim--Skolem proof.
  \item \emph{Syntax counting} (\cref{cor:syntax-counting}). Proved as
        \cref{thm:p01-w02-p01-011-syntax-counting} from clauses~2 and~4.
\end{enumerate}
The Cantor--Schr\"oder--Bernstein theorem, which the upward proof uses to
pass from two inequalities to an equality, is proved where it is stated, in
\cref{sec:set-prerequisites-05}.

\subsection{The bound \texorpdfstring{\(\abs X\leq\mu\)}{|X| <= mu} and the cardinal \texorpdfstring{\(\mu\)}{mu}}
\label{sec:p02-w02b-p02-029-bound}

\begin{hblemma}{thm:p02-w02b-p02-029-bounds}{proved=annex:p02-w02b-p02-029-bounds}{Closure of the bound \(\abs X\leq\mu\)}
Let \(\mu\) be an infinite cardinal and let \(X\), \(Y\) be sets with
\(\abs X\leq\mu\) and \(\abs Y\leq\mu\).
\begin{enumerate}
  \item Every \(Z\subseteq X\) satisfies \(\abs Z\leq\mu\).
  \item Every countable set, and in particular every finite set, \(Z\)
        satisfies \(\abs Z\leq\mu\).
  \item \(\abs{X\cup Y}\leq\mu\) and \(\abs{X\times Y}\leq\mu\).
  \item For every function \(g\) with domain \(X\), the image
        \(g[X]\) satisfies \(\abs{g[X]}\leq\mu\).
  \item Under the Axiom of Choice, \(\abs{X^{<\omega}}\leq\mu\).
\end{enumerate}
Items~1--4 use no form of the Axiom of Choice.
\end{hblemma}

Item~4 is the step that the inherited proofs pass over most quickly: a set
of values \(f^{\mathcal M}(\bar a)\) or \(F_\varphi(\bar a)\) is the image of
a set of pairs of a symbol and a tuple, and its size is bounded by the size
of that set of pairs. The proof chooses, for each value, the least code of a
pair that produces it, so it needs no choice.

\begin{hblemma}{thm:p02-w02b-p02-029-mu}{proved=annex:p02-w02b-p02-029-mu}{The cardinal of the downward bound}
Assume the Axiom of Choice. Let \(\Lang=(\mathcal F,\mathcal R,\ldots)\) be a
signature and \(A\) a set. Then the cardinals \(\abs A\) and
\(\abs\Lang=\abs{\mathcal F\cup\mathcal R}\) exist, the cardinal
\[
  \mu=\max(\abs A,\abs\Lang,\aleph_0)
\]
exists and is an infinite cardinal, and \(\abs A\leq\mu\),
\(\abs{\mathcal F\cup\mathcal R}\leq\mu\) and \(\abs\N\leq\mu\) hold in the
sense of injections.
\end{hblemma}

\begin{hblemma}{thm:p02-w02b-p02-029-finite-injections}{proved=annex:p02-w02b-p02-029-finite-injections}{Finite pieces of an infinite set}
Let \(D\) be an infinite set, that is, a set bijective with no \(n\in\N\).
For every \(n\in\N\) there is an injection \(n\to D\), and for every finite
set \(I\) there is an injection \(I\to D\). No form of the Axiom of Choice
is used.
\end{hblemma}

\subsection{Witness functions from one well-order}
\label{sec:p02-w02b-p02-029-witnesses}

The downward proof chooses a witness function \(F_\varphi\) for each formula
\(\exists y\,\varphi(y,\bar x)\). The notation \(\varphi(y,x_1,\ldots,x_k)\)
displays a list of variables, so the index of such a function is a formula
together with that list.

\begin{hbdefinition}{def:p02-w02b-p02-029-patterns}{Existential pattern}
Let \(\Lang\) be a signature. An \emph{existential pattern} of \(\Lang\) is
a pair \(e=(\varphi,(y,x_1,\ldots,x_k))\) of a formula
\(\varphi\in\mathsf{Fm}(\Lang)\) and a finite sequence of
\(k+1\) pairwise distinct variables, \(k\in\N\), such that every free
variable of \(\varphi\) is among \(y,x_1,\ldots,x_k\); \(k\) is the
\emph{arity} of \(e\). The set of existential patterns of \(\Lang\) is
written \(\mathcal E(\Lang)\). For a structure \(\mathcal M\),
\(b\in\dom\mathcal M\) and \(\bar a=(a_1,\ldots,a_k)\in(\dom\mathcal M)^k\),
\(\mathcal M\entails\varphi(b,\bar a)\) means that \(\varphi\) is satisfied
under an assignment sending \(y\) to \(b\) and \(x_i\) to \(a_i\); by the
coincidence of satisfaction for assignments agreeing on the free variables
this does not depend on the values of other variables. Then
\(\mathcal M\entails\exists y\,\varphi(y,\bar a)\) holds iff
\(\mathcal M\entails\varphi(b,\bar a)\) for some \(b\in\dom\mathcal M\).
\end{hbdefinition}

\begin{hblemma}{thm:p02-w02b-p02-029-skolem-functions}{proved=annex:p02-w02b-p02-029-skolem-functions}{Witness functions}
Assume the Axiom of Choice. Let \(\mathcal M\) be an \(\Lang\)-structure and
\(m_0\in\dom\mathcal M\). There is a function \(e\mapsto F_e\) on
\(\mathcal E(\Lang)\) such that, for every pattern
\(e=(\varphi,(y,x_1,\ldots,x_k))\), \(F_e\) is a function
\((\dom\mathcal M)^k\to\dom\mathcal M\) and, for every
\(\bar a\in(\dom\mathcal M)^k\),
\begin{enumerate}
  \item if \(\mathcal M\entails\exists y\,\varphi(y,\bar a)\), then
        \(\mathcal M\entails\varphi(F_e(\bar a),\bar a)\);
  \item otherwise \(F_e(\bar a)=m_0\).
\end{enumerate}
The Axiom of Choice is used once, through the well-ordering theorem applied
to \(\dom\mathcal M\); after that the whole family is defined outright.
\end{hblemma}

In \cref{thm:downward-ls} the function written \(F_\varphi\) is \(F_e\) for
the pattern \(e\) of the displayed formula \(\varphi(y,\bar x)\).

\subsection{The stages of the downward construction}
\label{sec:p02-w02b-p02-029-stage-bounds}

\begin{hbdefinition}{def:p02-w02b-p02-029-stages}{Stages}
Let \(\mathcal M\), \(A\subseteq\dom\mathcal M\), \(m_0\) and
\((F_e)_{e\in\mathcal E(\Lang)}\) be as in
\cref{thm:p02-w02b-p02-029-skolem-functions}. For
\(X\subseteq\dom\mathcal M\) put
\[
\begin{aligned}
  \operatorname{Fn}(X)&=\{f^{\mathcal M}(\bar a):f\in\mathcal F
      \text{ of arity }k\geq1,\ \bar a\in X^k\},\\
  \operatorname{Wt}(X)&=\{F_e(\bar a):e\in\mathcal E(\Lang)
      \text{ of arity }k,\ \bar a\in X^k\},\\
  \operatorname{St}(X)&=X\cup\operatorname{Fn}(X)\cup\operatorname{Wt}(X).
\end{aligned}
\]
Let \(A_0=A\cup\{m_0\}\cup\{c^{\mathcal M}:c\in\mathcal F\text{ of arity
}0\}\), and let \((A_n)_{n\in\N}\) be the unique function on \(\omega=\N\)
with \(A_{n+1}=\operatorname{St}(A_n)\), given by transfinite recursion
(\cref{thm:p01-w02-p01-013-recursion}) with \(\delta=\omega\). Put
\(N=\bigcup_{n\in\N}A_n\).
\end{hbdefinition}

These are exactly the sets \(A_n\) and \(N\) of the proof of
\cref{thm:downward-ls}.

\begin{hbtheorem}{thm:p02-w02b-p02-029-stages}{proved=annex:p02-w02b-p02-029-stages}{Stage bounds and closure}
Assume the Axiom of Choice, and let
\(\mu=\max(\abs A,\abs\Lang,\aleph_0)\) and the stages be as in
\cref{def:p02-w02b-p02-029-stages}.
\begin{enumerate}
  \item If \(X\subseteq\dom\mathcal M\) and \(\abs X\leq\mu\), then
        \(\abs{\operatorname{St}(X)}\leq\mu\).
  \item \(\abs{A_n}\leq\mu\) for every \(n\in\N\), and
        \(\abs N\leq\mu\).
  \item \(A_m\subseteq A_n\) whenever \(m\leq n\); \(A\subseteq N\); and
        \(N\neq\varnothing\).
  \item For every \(k\in\N\) and \(\bar a\in N^k\) there is \(n\in\N\)
        with \(\bar a\in A_n^k\).
  \item \(N\) contains \(c^{\mathcal M}\) for every constant symbol \(c\) and
        is closed under \(f^{\mathcal M}\) for every function symbol \(f\) of
        positive arity.
  \item For every pattern \(e=(\varphi,(y,x_1,\ldots,x_k))\) and
        \(\bar a\in N^k\): if
        \(\mathcal M\entails\exists y\,\varphi(y,\bar a)\), then
        \(\mathcal M\entails\varphi(b,\bar a)\) for some \(b\in N\).
\end{enumerate}
\end{hbtheorem}

Items~2, 3 and~5 supply the steps of the proof of \cref{thm:downward-ls}
that bound every stage and the union by \(\mu\) and that make the union a
nonempty superset of \(A\) closed under the interpreted function symbols;
items~4 and~6 are its final paragraph, which with the Tarski--Vaught test
gives \(\mathcal N\preccurlyeq\mathcal M\) for the induced substructure
\(\mathcal N\) on \(N\).

\subsection{The cardinal steps of the upward theorem}
\label{sec:p02-w02b-p02-029-upward}

\begin{hblemma}{thm:p02-w02b-p02-029-fresh}{proved=annex:p02-w02b-p02-029-fresh}{Fresh constant symbols}
Let \(\Gamma\) be a set and \(\kappa\) an ordinal. There are objects
\(c_\alpha\), \(\alpha<\kappa\), pairwise distinct and none a member of
\(\Gamma\). No form of the Axiom of Choice is used.
\end{hblemma}

Applied to the alphabet \(\Gamma=\mathcal F\cup\mathcal R\cup\mathsf{Var}\cup
T_{\mathrm{log}}\), with \(T_{\mathrm{log}}\) the reserved logical and
punctuation tokens of \cref{sec:fol-syntax-semantics-01}, the lemma supplies
the new constants of the upward proof. The signature \(\Lang^*\) obtained by
adding them as constant symbols then satisfies the disjointness requirement
of a signature.

\begin{hbtheorem}{thm:p02-w02b-p02-029-upward-count}{proved=annex:p02-w02b-p02-029-upward-count}{Cardinal steps of upward L\"owenheim--Skolem}
Assume the Axiom of Choice. Let \(\Lang\) be a signature, \(\kappa\) a
cardinal with \(\kappa\geq\max(\abs\Lang,\aleph_0)\), \(c_\alpha\)
(\(\alpha<\kappa\)) new constant symbols as in
\cref{thm:p02-w02b-p02-029-fresh}, and \(\Lang^*\) the expanded signature.
\begin{enumerate}
  \item \(\abs{\Lang^*}=\kappa\).
  \item If \(\Delta\) is a finite set of \(\Lang^*\)-formulas, the set
        \(I_\Delta\) of \(\alpha<\kappa\) such that \(c_\alpha\) occurs in a
        member of \(\Delta\) is finite, and every infinite set \(D\) admits
        an injection \(I_\Delta\to D\).
  \item If \(\mathcal M^*\) is an \(\Lang^*\)-structure satisfying
        \(c_\alpha\ne c_\beta\) for all \(\alpha<\beta<\kappa\), then
        \(\alpha\mapsto c_\alpha^{\mathcal M^*}\) is a bijection from
        \(\kappa\) onto \(A=\{c_\alpha^{\mathcal M^*}:\alpha<\kappa\}\), so
        \(\abs A=\kappa\), and
        \(\max(\abs A,\abs{\Lang^*},\aleph_0)=\kappa\).
  \item If, in addition, \(A\subseteq N\) and \(\abs N\leq\kappa\), then
        \(\abs N=\kappa\).
\end{enumerate}
\end{hbtheorem}

Item~1 makes downward L\"owenheim--Skolem in \(\Lang^*\) give the bound
\(\kappa\) of item~3; item~2 is the step ``assign distinct elements to the
finitely many constants involved''; items~3 and~4 are the steps ``\(A\) has
cardinality \(\kappa\)'' and ``containment of \(A\) gives cardinality at
least \(\kappa\), so the cardinality is \(\kappa\)''.

\begin{hbmethod}{def:p02-w02b-p02-029-method}{Checking the set theory of a L\"owenheim--Skolem argument}
To verify the cardinal bookkeeping in a downward argument of the form of
\cref{thm:downward-ls}:
\begin{enumerate}
  \item Fix \(\mu=\max(\abs A,\abs\Lang,\aleph_0)\)
        (\cref{thm:p02-w02b-p02-029-mu}); the Axiom of Choice is needed for
        the cardinals \(\abs A\), \(\abs\Lang\) and their comparison.
  \item Fix one well-order of the domain and take least witnesses
        (\cref{thm:p02-w02b-p02-029-skolem-functions}). Side condition: the
        index of a witness function is a pattern, a formula together with
        its displayed variables.
  \item Present each new stage as the image of a set of pairs (symbol or
        pattern, finite tuple) and bound that set with products and finite
        sequences (\cref{thm:p02-w02b-p02-029-bounds}, items~3--5).
  \item Bound the union of the countably many stages with clause~4 of the
        finite-coding theorem; this needs the Axiom of Choice because the
        induction gives only the existence of an injection for each stage.
  \item Check that every finite tuple from the union lies in one stage
        (\cref{thm:p02-w02b-p02-029-stages}, item~4) before applying the
        Tarski--Vaught test.
\end{enumerate}
For an upward argument, add: constants are fresh
(\cref{thm:p02-w02b-p02-029-fresh}); the expanded signature still has
cardinality \(\kappa\); and the final equality comes from two injections and
the Cantor--Schr\"oder--Bernstein theorem
(\cref{thm:p02-w02b-p02-029-upward-count}).
\end{hbmethod}

\begin{hbexample}{def:p02-w02b-p02-029-example}{A countable elementary substructure of the subsets of \(\N\)}
Let \(\Lang\) have a single binary relation symbol \(\sqsubseteq\) and no
function symbols, and let \(\mathcal M\) have domain \(\mathcal P(\N)\) and
interpret \(\sqsubseteq\) as inclusion; let \(A=\{\varnothing\}\). Then
\(\mathcal F\cup\mathcal R=\{\sqsubseteq\}\), so \(\abs\Lang=1\);
\(\abs A=1\); and \(\mu=\max(1,1,\aleph_0)=\aleph_0\) by item~5 of
\cref{thm:p01-w02-p01-008-cardinal-assignment}. The alphabet adds the
countable set of variables and the finitely many logical tokens to
\(\{\sqsubseteq\}\), so \(\mathsf{Fm}(\Lang)\) is countable by syntax
counting, and so is \(\mathcal E(\Lang)\)
(\cref{thm:p02-w02b-p02-029-bounds}, items~3 and~5, with \(\mu=\aleph_0\)).
With \(m_0=\varnothing\), the stage \(A_0=\{\varnothing\}\) is finite. Since
there are no function symbols, \(\operatorname{Fn}(X)=\varnothing\), and
each later stage adds only the witnesses \(F_e(\bar a)\) for tuples
\(\bar a\) from the previous stage: for instance the pattern
\((\neg\,y\sqsubseteq x_1,(y,x_1))\) at \(x_1\mapsto\varnothing\) contributes
a nonempty subset of \(\N\) to \(A_1\), and the pattern
\((\forall x_2\,y\sqsubseteq x_2,(y))\) of arity \(0\) contributes
\(\varnothing\) again. The witnesses are least elements in a well-order of
\(\mathcal P(\N)\) supplied by the well-ordering theorem, which describes no
particular one, so the members of \(N\) other than \(\varnothing\) are not
determined; but \cref{thm:p02-w02b-p02-029-stages} gives
\(\abs N\leq\aleph_0\) whatever the well-order is, and \cref{thm:downward-ls}
gives an elementary substructure with domain \(N\). Cantor's theorem in
\cref{sec:set-prerequisites-05} gives \(\abs\N<\abs{\mathcal P(\N)}\), that
is, no bijection \(\N\to\mathcal P(\N)\) exists; since
\(n\mapsto\{n\}\) is an injection \(\N\to\mathcal P(\N)\), the
Cantor--Schr\"oder--Bernstein theorem excludes an injection
\(\mathcal P(\N)\to\N\), so \(N\neq\mathcal P(\N)\) and the elementary
substructure is proper. It is infinite by item~2 of
\cref{thm:p02-w02b-p02-029-terms-needed}, so \(\abs N=\aleph_0\): the bound
\(\mu\) is attained through the term \(\aleph_0\) alone, although \(\abs A\)
and \(\abs\Lang\) are both \(1\).
\end{hbexample}

\begin{hblemma}{thm:p02-w02b-p02-029-terms-needed}{proved=annex:p02-w02b-p02-029-terms-needed}{Each term of the downward bound is needed}
Assume the Axiom of Choice. Let \(\mathcal M\) be an \(\Lang\)-structure,
\(A\subseteq\dom\mathcal M\), and \(\mathcal N\preccurlyeq\mathcal M\) with
\(A\subseteq\dom\mathcal N\).
\begin{enumerate}
  \item \(\abs A\leq\abs{\dom\mathcal N}\).
  \item If \(\dom\mathcal M\) is infinite, then \(\dom\mathcal N\) is
        infinite and \(\aleph_0\leq\abs{\dom\mathcal N}\).
  \item \(\abs{C^{\mathcal M}}\leq\abs{\dom\mathcal N}\), where
        \(C^{\mathcal M}=\{c^{\mathcal M}:c\in\mathcal F\text{ of arity }0\}\).
  \item For every infinite cardinal \(\kappa\) there are a signature
        \(\Lang_\kappa\) with \(\abs{\Lang_\kappa}=\kappa\) and an
        \(\Lang_\kappa\)-structure \(\mathcal M_\kappa\) with domain
        \(\kappa\) such that every substructure of \(\mathcal M_\kappa\), and
        so every elementary substructure, has domain \(\kappa\).
\end{enumerate}
\end{hblemma}

\begin{hbwarning}{def:p02-w02b-p02-029-limits}{What the prerequisites do and do not provide}
\begin{enumerate}
  \item \emph{No term of the maximum can be deleted.} By item~1 of
        \cref{thm:p02-w02b-p02-029-terms-needed}, a bound without \(\abs A\)
        fails whenever \(A\) is larger than the other two terms. By item~2,
        \(\aleph_0\) cannot be deleted: in
        \cref{def:p02-w02b-p02-029-example}, \(\abs A=\abs\Lang=1\) but every
        elementary substructure is infinite. By item~4, \(\abs\Lang\) cannot
        be deleted: for \(\kappa>\aleph_0\) and \(A=\varnothing\),
        \(\mathcal M_\kappa\) has no elementary substructure of cardinality
        below \(\kappa\), while \(\max(\abs A,\aleph_0)=\aleph_0\).
  \item \emph{Where the Axiom of Choice enters.} In the downward proof:
        the cardinals \(\abs A\), \(\abs\Lang\) and their maximum
        (\cref{thm:p02-w02b-p02-029-mu}, through cardinal assignment and the
        well-ordering theorem); the well-order of the domain
        (\cref{thm:p02-w02b-p02-029-skolem-functions}); finite sequences
        (item~5 of \cref{thm:p02-w02b-p02-029-bounds}, through clause~2 of
        the finite-coding theorem); and the countable union of the stages
        (clause~4). In the upward proof, additionally, the cardinal
        assignment of item~3 of \cref{thm:p02-w02b-p02-029-upward-count}.
        These remarks record what the proofs here use. They do not claim that
        either theorem fails without the Axiom of Choice; that question is
        not settled by anything proved in this book.
  \item \emph{Existence only.} The elementary substructure depends on the
        well-order chosen for the domain; nothing here describes its members
        beyond \(A\), the constants and their closure
        (\cref{def:p02-w02b-p02-029-example}).
  \item \emph{The upward hypothesis \(\kappa\geq\abs\Lang\).} Item~1 of
        \cref{thm:p02-w02b-p02-029-upward-count} uses it to keep
        \(\abs{\Lang^*}=\kappa\); \cref{thm:p02-w02-p02-033-hypotheses-needed}
        shows that it cannot be dropped from the upward theorem.
\end{enumerate}
\end{hbwarning}

\subsection{Relationships and sources}
\label{sec:p02-w02b-p02-029-sources}

The statements of \cref{thm:downward-ls,thm:upward-ls} and their proofs are
those of the frozen reference \cite[lines 5043--5131, section ``The
L\"owenheim--Skolem theorems'']{FrozenFormalLogicCompactnessApplications},
and the four prerequisites are stated there at
\cite[lines 657--660 (well-ordering theorem), 662--667 (cardinal
assignment), 679--694 (finite coding) and 696 onward (syntax
counting)]{FrozenFormalLogicFiniteCoding}, where the first three carry a
citation in place of a proof. The proofs of the cited
results are those of Part~1 listed in \cref{sec:p02-w02b-p02-029-where}; the
lemmas and theorems of this section, and their proofs, are editorial
completions built from the definitions of the frozen reference and are not
attributed to any external presentation. The same bookkeeping serves the
version of upward L\"owenheim--Skolem through elementary diagrams
(\cref{thm:p02-upward-ls-elementary-extension}) and its inclusion form
(\cref{thm:p02-w02-p02-033-upward-ls-inclusion}), whose proofs apply
\cref{thm:downward-ls} in an expanded signature of cardinality \(\kappa\).
